% ===========================================================================
\section{Implementation Plan}
% The plan in the Evaluation 1 report, against what was delivered.
% ===========================================================================

The Evaluation~1 report planned the work in phases: tier 2 working end to end
(M1), then tier 3 on Sepolia (M2), a feature freeze with auditing and usability
testing, and measurement. Table~\ref{tab:plan-executed} sets each against what
was delivered.

\begin{table}[!htb]
  \centering
  \caption{The plan against what was delivered.}
  \label{tab:plan-executed}
  \small
  \begin{tabular}{@{}L{2.3cm}L{6.2cm}L{7.3cm}@{}}
    \toprule
    \textbf{Phase} & \textbf{Planned}                                             & \textbf{Delivered}                                                                                            \\
    \midrule
    1 (M1)         & Tier 2 end to end on a local node; first gas snapshot          & The lifecycle, factory, registrar, indexer, admin console and booth, with a browser test and the tier 2 gas snapshot \\
    2 (M2)         & Tier 3 anonymous and gasless voting on Sepolia                 & Built and deployed on Sepolia, where complete anonymous elections run end to end in real time \\
    3 (freeze)     & Approval voting if time allowed, Slither, reorgs, audit with the backend off, usability test & Approval voting, Slither, coverage, reorg handling, and the audit panel, which passes with the API and relayer unreachable. The usability test was run as a simulation \\
    4 (measure)    & Every parameter measured                                       & A1, B1, B2 on a laptop and an Android phone, C1 to C3, D1 to D5, E1, E2, and E3 by simulation (Section~6) \\
    \bottomrule
  \end{tabular}
\end{table}

Whether phones could prove ballots was settled by measurement: proving
a ballot takes 0.57\,s at two voters and 1.52\,s at 65,536 on one ARM64 CPU
core, and in the browser on an Android phone under two seconds up to 4,096
voters and 9.1\,s at 65,536 (B2). That is fast enough that proving stays on the
voter's device and the server-side fallback is not needed.

Some items were built differently from the plan. The design changes, such as
the relayer calling \texttt{castVote} itself without a forwarder, documents on
our own IPFS node with copies on Pinata, and gas reported under both mainnet's
and Glamsterdam's rules, are set out with their reasons in
Section~\ref{sec:req-changes} and Table~\ref{tab:divergences}. The proving files
are also checked differently: the plan was to compare them with hashes
published in the Semaphore repository, and they are instead pinned by SHA-256
against the npm package of the same release, whose own integrity hash matched
the registry. Every file the booth downloads is checked against those pins.

Approval voting, optional in the plan, is built and fixed per election when it
is created. Two things were added that the plan did not have: Google sign-in
for the institute's domain in place of an emailed code, and a QR code or
private link that carries the voting key to a second device.

\subsection{Contingencies}

The Evaluation~1 plan named five risks. Three of its responses were built as
planned, and the two fallbacks were never needed:

\begin{itemize}
  \item \emph{Relayer unavailable.} Any wallet can submit an anonymous
        ballot. The booth always offers this, and recommends it when the relayer
        reports a failure or voting closes within five minutes.
  \item \emph{Relayer funds drained.} The relayer pays only for elections run
        by the deployment's authorities, since anyone can create one. It
        verifies the proof, checks that the nullifier is unspent and not already
        queued, and simulates the call before it pays, so it pays for at most
        one counted ballot per registered voter. A ballot that loses a race, for example to the close
        of voting, still costs gas.
  \item \emph{Not enough Sepolia ETH.} Gas and scaling measurements run on a
        local node, as planned. Test ETH is not a constraint for the rest:
        deploying both tiers cost 0.00061\,ETH (Table~\ref{tab:deployment}).
  \item \emph{Proving moves to a server.} Not needed, as above.
  \item \emph{Tier 3 not ready, so tier 2 is demonstrated.} Not needed: tier 3
        runs on Sepolia.
\end{itemize}

\subsection{Future work}

Six pieces of work would matter most beyond this project:

\begin{itemize}
  \item \textbf{Coercion resistance.} A voter can still show a buyer how they
        voted, for example by voting in front of them, and the key backup is one
        more thing that can be handed over. A coercion-resistant design makes
        it infeasible for a coercer to tell whether the voter
        complied~\cite{jcj}. On Ethereum, MACI~\cite{maci} lets a voter cast a
        later vote that nullifies an earlier one, so a vote cast under pressure
        need not stand. The price is a coordinator who decrypts the votes and must
        be trusted for collusion resistance, though not for the count. A tier
        built on MACI's model would be the natural next step.
  \item \textbf{Tally secrecy while voting is open.} In \projectname{} the running count is
        public, which can sway late voters. In Helios~\cite{helios}, the
        bulletin board shows each ballot only in encrypted form, and the
        ballots are shuffled and decrypted once the election closes; but whoever holds the decryption
        key must be trusted not to use it early. In the Open Vote Network on
        Ethereum~\cite{openvote}, nobody holds such a key, but the final tally
        can be computed only after the last registered voter has voted, and a
        voter who abstains blocks it. Its authors designed it for small
        boardroom elections, not a university electorate.
  \item \textbf{Production identity.} The registrar accepts Google sign-in
        restricted to the institute's Google domain~\cite{googleidtoken}. An
        institute deployment would instead connect to the institute's own
        single sign-on, and publish which attribute marks an eligible voter.
        The registrar would remain the only trusted party for eligibility.
  \item \textbf{A relayer that cannot censor selectively.} One relayer sees
        when ballots arrive, from which network address, and could quietly
        drop some. Several independent relayers, with the booth spreading
        ballots between them and falling back to direct submission, would
        remove any single point of censorship.
  \item \textbf{Key portability.} A voting key lives in one browser. The QR code
        and the backup file move it, but only if the voter thinks to. A key
        derived from a passkey that syncs between the voter's own devices would
        remove the need to move it.
  \item \textbf{Pages a voter can check.} The voter trusts whoever serves the
        web app not to send the key elsewhere. A build that anyone can
        reproduce from the source, published by its content hash on IPFS,
        would let a voter, or someone they trust, check that the pages they
        load are the published code.
\end{itemize}
